\clearpage\gdef\otChapNum{0}\pagestyle{otFront}\phantomsection\label{defis-debut}\pdfbookmark[0]{Défis d’approfondissement}{bm-defis}
\otPageTitre{Défis d’approfondissement}{Un parcours facultatif pour aller plus loin}
Les 16 défis demandent davantage d’initiative que les exercices de chapitre. Ils restent fondés sur des outils accessibles à partir de la Seconde ; les prolongements sont définis dans l’énoncé. Les pistes d’approfondissement du programme sont toutes repérées dans la matrice finale.\par\medskip
Chercher d’abord sans aide. Les indications sont séparées des énoncés ; ouvrir la première, puis la seconde seulement si nécessaire. Prévoir de 30 à 90 minutes selon le défi, éventuellement en plusieurs séances.\par\medskip
\begin{savoirfaire}[Conserver une trace de recherche]Dessiner, tester un cas, écrire une conjecture, puis revenir aux hypothèses. Dans le corrigé, repérer l’idée qui débloque le problème et la formulation de la conclusion.\end{savoirfaire}
\begin{tabularx}{\linewidth}{@{}lY@{}}\toprule
D01 à D03 & Nombres, identités et inégalités.\\[2mm]
D04 à D11 & Barycentres, concours de droites, aires et lieux.\\[2mm]
D12 à D13 & Paramètres et dispersion.\\[2mm]
D14 à D16 & Tableaux, contrôles et dénombrement.\\\bottomrule
\end{tabularx}
\editionPage{defi-ex-1}{D01 · Décimales finies et périodes}{Recherche et rédaction}{Chapitres 1 et 2 · \hyperref[defi-cor-1-1]{Corrigé p.\,\pageref*{defi-cor-1-1}}}
\begin{blocEdition}{defi-ex-1}
On étudie les écritures décimales de fractions irréductibles positives $p/q$.
\begin{questions}
\item Montrer que $p/q$ est décimal si et seulement si $q$ ne contient que des facteurs premiers 2 et 5. On peut utiliser l’unicité de la décomposition en facteurs premiers.
\item Effectuer la division de $1$ par 7 jusqu’au retour du premier reste. Donner la période.
\item Expliquer pourquoi la division d’un rationnel donne une écriture finie ou une écriture périodique à partir d’un certain rang.
\end{questions}
\end{blocEdition}
\editionPage{defi-ex-2}{D02 · Deux identités à construire}{Recherche et rédaction}{Chapitre 3 · \hyperref[defi-cor-2-1]{Corrigé p.\,\pageref*{defi-cor-2-1}}}
\begin{blocEdition}{defi-ex-2}
\begin{questions}
\item Développer $(a+b+c)^2$ en faisant apparaître les carrés et les doubles produits.
\item Déduire le développement de $(a+b)^3$ en multipliant $(a+b)^2$ par $(a+b)$.
\item Utiliser ces identités pour calculer $103^3$ et $(10+2+1)^2$.
\end{questions}
\end{blocEdition}
\editionPage{defi-ex-3}{D03 · Moyennes et rectangle optimal}{Recherche et rédaction}{Chapitres 3 et 7 · \hyperref[defi-cor-3-1]{Corrigé p.\,\pageref*{defi-cor-3-1}}}
\begin{blocEdition}{defi-ex-3}
Pour $a,b>0$, on appelle moyenne arithmétique $(a+b)/2$ et moyenne géométrique $\sqrt{ab}$.\begin{questions}
\item Partir de $(\sqrt a-\sqrt b)^2\ge0$ pour comparer ces deux moyennes et préciser le cas d’égalité.
\item Deux côtés d’un rectangle positif ont une somme égale à 18. Déterminer son aire maximale.
\item À aire fixée à 36, déterminer le périmètre minimal.
\end{questions}\label{fig-defi-ex-3}
\figureGeo{defi-moyennes}

\end{blocEdition}
\editionPage{defi-ex-4}{D04 · Barycentres de deux et trois points}{Recherche et rédaction}{Chapitres 8 et 9 · \hyperref[defi-cor-4-1]{Corrigé p.\,\pageref*{defi-cor-4-1}}}
\begin{blocEdition}{defi-ex-4}
Pour des points de position $A$ et $B$ et des poids réels $\alpha,\beta$ de somme non nulle, on définit ici $G$ par $\alpha\vect{GA}+\beta\vect{GB}=\vec0$.\begin{questions}
\item Exprimer les coordonnées de $G$ en fonction des poids et des coordonnées des deux points.
\item Pour $A(0;0)$, $B(6;0)$ et les poids 2 et 1, construire $G$. Que devient la position avec les poids 2 et $-1$ ?
\item Étendre la formule à trois points de poids $\alpha,\beta,\gamma$ de somme non nulle. Application : $A(0;0)$, $B(6;0)$, $C(0;6)$, poids 1, 1 et 2.
\end{questions}\label{fig-defi-ex-4}
\figureGeo{defi-barycentres}

\end{blocEdition}
\editionPage{defi-ex-5}{D05 · Pourquoi les médianes se coupent}{Recherche et rédaction}{Chapitres 8 et 9 · \hyperref[defi-cor-5-1]{Corrigé p.\,\pageref*{defi-cor-5-1}}}
\begin{blocEdition}{defi-ex-5}
Soient trois points non alignés $A,B,C$. On note $G=(A+B+C)/3$ en coordonnées et $I$ le milieu de $[BC]$.\begin{questions}
\item Montrer $\vect{AG}=\dfrac23\vect{AI}$.
\item En déduire que les trois médianes sont concourantes en $G$, l’isobarycentre des trois points.
\item Pour $A(0;0)$, $B(6;0)$ et $C(0;6)$, calculer $G$ et vérifier les trois relations.
\end{questions}\label{fig-defi-ex-5}
\figureGeo{defi-medianes}

\end{blocEdition}
\editionPage{defi-ex-6}{D06 · Concours des hauteurs}{Recherche et rédaction}{Chapitre 9 ; parallèles et médiatrices · \hyperref[defi-cor-6-1]{Corrigé p.\,\pageref*{defi-cor-6-1}}}
\begin{blocEdition}{defi-ex-6}
Dans un triangle non aplati $ABC$, on construit $A'=B+C-A$, $B'=A+C-B$ et $C'=A+B-C$ en coordonnées. Ces points forment un grand triangle dont les côtés sont parallèles à ceux de $ABC$.\begin{questions}
\item Montrer que $A$ est le milieu de $[B'C']$, et faire les deux autres calculs analogues.
\item Montrer que $(B'C')$ est parallèle à $(BC)$.
\item En admettant ou en utilisant la preuve du concours des médiatrices du défi D08, en déduire que les trois hauteurs de $ABC$ sont concourantes.
\end{questions}\label{fig-defi-ex-6}
\figureGeo{defi-hauteurs}

\end{blocEdition}
\editionPage{defi-ex-7}{D07 · Une aire à partir de deux côtés et d’un angle}{Recherche et rédaction}{Chapitre 9 ; trigonométrie du triangle rectangle · \hyperref[defi-cor-7-1]{Corrigé p.\,\pageref*{defi-cor-7-1}}}
\begin{blocEdition}{defi-ex-7}
Dans un triangle, les côtés $CA=b$ et $CB=a$ forment l’angle $C$.\begin{questions}
\item Pour $C$ aigu, utiliser une hauteur pour établir $\mathcal A=\dfrac12ab\sin C$.
\item Pour $C$ obtus, on définit dans ce prolongement $\sin C=\sin(180^\circ-C)$. Reprendre la démonstration avec la hauteur extérieure au segment.
\item Calculer l’aire pour $a=8$, $b=5$, $C=30^\circ$, puis pour $C=150^\circ$. Les deux triangles ont-ils nécessairement le même troisième côté ?
\end{questions}\label{fig-defi-ex-7}
\figureGeo{defi-aire-angle}

\end{blocEdition}
\editionPage{defi-ex-8}{D08 · Concours des médiatrices et cercle circonscrit}{Recherche et rédaction}{Chapitre 9 · \hyperref[defi-cor-8-1]{Corrigé p.\,\pageref*{defi-cor-8-1}}}
\begin{blocEdition}{defi-ex-8}
Dans un triangle non aplati $ABC$, les médiatrices de $[AB]$ et $[AC]$ se coupent en $O$.\begin{questions}
\item Justifier qu’elles ne sont pas parallèles.
\item Montrer que $O$ appartient aussi à la médiatrice de $[BC]$.
\item En déduire l’existence et l’unicité du cercle passant par les trois sommets. Illustrer pour $A(0;0)$, $B(6;0)$, $C(2;4)$.
\end{questions}\label{fig-defi-ex-8}
\figureGeo{defi-circonscrit}

\end{blocEdition}
\editionPage{defi-ex-9}{D09 · Un lieu parabolique}{Recherche et rédaction}{Chapitres 9 et 10 · \hyperref[defi-cor-9-1]{Corrigé p.\,\pageref*{defi-cor-9-1}}}
\begin{blocEdition}{defi-ex-9}
On cherche les points $M(x;y)$ équidistants du point $F(0;2)$ et de l’axe des abscisses. La distance de $M$ à cet axe est $|y|$.\begin{questions}
\item Traduire l’équidistance par une égalité de carrés.
\item Montrer que le lieu a pour équation $y=1+x^2/4$.
\item Justifier la réciproque et déterminer son point le plus bas. Trouver les points de ce lieu sur $y=5$.
\end{questions}\label{fig-defi-ex-9}
\figureGeo{defi-parabole}

\end{blocEdition}
\editionPage{defi-ex-10}{D10 · Des inégalités décrivent une région}{Recherche et rédaction}{Chapitres 4 et 10 · \hyperref[defi-cor-10-1]{Corrigé p.\,\pageref*{defi-cor-10-1}}}
\begin{blocEdition}{defi-ex-10}
On cherche les couples $(x;y)$ tels que $x\ge0$, $y\ge0$, $x+y\le8$ et $x+2y\le12$.\begin{questions}
\item Représenter la région et déterminer ses sommets.
\item On souhaite maximiser $R=3x+4y$. Démontrer directement que $R\le28$ en combinant deux contraintes.
\item Trouver le cas d’égalité. Pourquoi comparer seulement quelques points dessinés n’aurait-il pas suffi à la preuve ?
\end{questions}\label{fig-defi-ex-10}
\figureGeo{defi-region}

\end{blocEdition}
\editionPage{defi-ex-11}{D11 · Carré et racine, deux courbes liées}{Recherche et rédaction}{Chapitres 5 et 6 · \hyperref[defi-cor-11-1]{Corrigé p.\,\pageref*{defi-cor-11-1}}}
\begin{blocEdition}{defi-ex-11}
\begin{questions}
\item Pour $a\ge0$, montrer que le point $(a;a^2)$ de la courbe carré correspond au point $(a^2;a)$ de la courbe racine.
\item Expliquer la symétrie par rapport à $y=x$, en précisant la restriction nécessaire de la courbe carré.
\item Déterminer les points d’intersection et comparer $\sqrt x$, $x$ et $x^2$ pour $0<x<1$, puis pour $x>1$.
\end{questions}\label{fig-defi-ex-11}
\figureGeo{defi-racine-carre}

\end{blocEdition}
\editionPage{defi-ex-12}{D12 · Un paramètre et deux extrémums}{Recherche et rédaction}{Chapitres 6 et 7 · \hyperref[defi-cor-12-1]{Corrigé p.\,\pageref*{defi-cor-12-1}}}
\begin{blocEdition}{defi-ex-12}
Pour $a\in\mathbb R$, on étudie $f_a(x)=(x-a)^2$ sur $[0;4]$.\begin{questions}
\item Discuter le minimum selon la position de $a$ par rapport à $[0;4]$.
\item Déterminer le maximum en comparant les distances aux extrémités.
\item Choisir $a$ pour que ce maximum soit le plus petit possible. Donner sa valeur et prouver l’unicité du choix.
\end{questions}
\end{blocEdition}
\editionPage{defi-ex-13}{D13 · Optimiser une dispersion}{Recherche et rédaction}{Chapitre 12 · \hyperref[defi-cor-13-1]{Corrigé p.\,\pageref*{defi-cor-13-1}}}
\begin{blocEdition}{defi-ex-13}
Une série comporte trois valeurs $0,t,6$, où $0\le t\le6$.\begin{questions}
\item Exprimer la moyenne et la médiane.
\item Montrer que la moyenne des carrés des écarts à la moyenne vaut $\dfrac29(t-3)^2+6$.
\item Pour quelle valeur de $t$ l’écart type est-il minimal ? Est-ce aussi le choix qui donne une moyenne de 3 ?
\end{questions}
\end{blocEdition}
\editionPage{defi-ex-14}{D14 · Des marges qui ne fixent pas l’intersection}{Recherche et rédaction}{Chapitres 13 et 14 · \hyperref[defi-cor-14-1]{Corrigé p.\,\pageref*{defi-cor-14-1}}}
\begin{blocEdition}{defi-ex-14}
Dans une population de 100 personnes, 60 pratiquent $A$ et 50 pratiquent $B$. On note $k$ le nombre qui pratiquent les deux.\begin{questions}
\item Compléter les quatre cases du tableau en fonction de $k$.
\item Déterminer toutes les valeurs entières possibles de $k$.
\item En déduire l’intervalle possible de la fréquence de $B$ parmi $A$. Les marges suffisent-elles à donner une seule valeur ?
\end{questions}
\end{blocEdition}
\editionPage{defi-ex-15}{D15 · Défauts rares et conditionnement inversé}{Recherche et rédaction}{Chapitres 13 et 14 ; prolongement facultatif · \hyperref[defi-cor-15-1]{Corrigé p.\,\pageref*{defi-cor-15-1}}}
\begin{blocEdition}{defi-ex-15}
Un contrôle industriel détecte 90\,\% des défauts et signale à tort 5\,\% des pièces conformes. La proportion de défauts est $p$, avec $0<p<1$. Ce défi prolonge le cours en calculant aussi la proportion globale de signaux positifs.\begin{questions}
\item Dans une population de taille $N$, exprimer les nombres théoriques de vrais et faux positifs en fonction de $p,N$.
\item Montrer que la proportion de défauts parmi les signaux positifs vaut $\dfrac{0{,}9p}{0{,}9p+0{,}05(1-p)}$.
\item Pour quelles valeurs de $p$ un signal positif correspond-il à un défaut avec une probabilité d’au moins $1/2$ ? Comparer $p=1\,\%$ et $p=10\,\%$.
\end{questions}\label{fig-defi-ex-15}
\figureGeo{defi-test}

\end{blocEdition}
\editionPage{defi-ex-16}{D16 · Trois dés, deux sommes}{Recherche et rédaction}{Chapitres 14 et 15 · \hyperref[defi-cor-16-1]{Corrigé p.\,\pageref*{defi-cor-16-1}}}
\begin{blocEdition}{defi-ex-16}
On lance trois dés équilibrés indépendants, distingués par leur couleur. Deux personnes pensent que les sommes 9 et 10 ont la même probabilité parce qu’elles se décrivent chacune par six triplets non ordonnés.\begin{questions}
\item Lister ces six triplets pour chaque somme, avec leurs termes dans l’ordre croissant.
\item Compter pour chaque triplet combien d’ordres différents il représente, puis comparer les probabilités.
\item Écrire un programme à trois boucles qui vérifie le dénombrement et expliquer le rôle des 216 issues.
\end{questions}
\end{blocEdition}
